% ==========================================================================
%  Section 2 -- Setting.  Owner: W02.  Frozen label: sec:setting.
%
%  This file fixes, and proves nothing: the torus and its test space
%  (def:torus), the pairing normalisation (rem:v-normalisation), the index
%  sets and the three weights (def:modes), the norm the weights are built
%  from (def:weight-supnorm), the negative-Sobolev carriers (def:carrier-hs),
%  the fixed-volume regime and the N/n dictionary (hyp:volume), the lattice
%  configuration space and the two symbols (eq:lattice-symbol,
%  eq:continuum-symbol), the two borrowed facts (ext:fourier-basis,
%  ext:summability), the dimensional caveat (rem:dimension) and the four
%  conventions a reader can get wrong (rem:conventions).
% ==========================================================================
\section{Setting}\label{sec:setting}

This section fixes the objects and the conventions that the rest of the paper
uses, and it proves nothing substantive. \Cref{def:torus} fixes the torus, its
test space and the Fourier coefficient; \cref{rem:v-normalisation} the
normalisation of the pairings; \cref{def:modes,def:weight-supnorm} the index
sets and the three weights; \cref{def:carrier-hs} the negative-Sobolev sequence
spaces that carry the measures and the energy. \Cref{hyp:volume} fixes the
fixed-volume regime and the dictionary between the two ways of indexing it. Two
background facts are stated: one is borrowed, as an external
theorem (\cref{ext:fourier-basis}), and the other is proved outright below
(\cref{ext:summability}). No probabilistic object is
introduced here: the lattice laws belong to \cref{sec:lattice}.

\subsection{The torus, its test space, and the Fourier coefficient}

\begin{definition}[The periodic torus and its test space]\label{def:torus}
Let $L>0$. The \emph{periodic torus} is $\Torus:=\R^2/(L\Z)^2$, identified
throughout with the half-open cell $[0,L)^2$; its volume is $V:=L^2$. A
\emph{test function} is a bounded continuous $f:\R^2\to\R$ that is
$L$-periodic in each coordinate,
\[
  f(x+Le_i)=f(x)\qquad\text{for } i=1,2 \text{ and all } x\in\R^2 ,
\]
where $e_1,e_2$ are the standard unit vectors, so a test function is a bounded
continuous real function on $\Torus$. The
trigonometric polynomials are uniformly dense in the set of test functions ---
a standard fact, recorded here for orientation; no step below spends it, because
every approximation in this paper is a finite-mode truncation of a series whose
terms are already explicit. The \emph{Fourier coefficient} of
a test function $f$ at the mode $k\in\Z^2$ is
\begin{equation}\label{eq:fourier-coefficient}
  F(k):=\int_{[0,L]^2} f(x)\,e^{-2\pi i\,k\cdot x/L}\,\dd x ,
\end{equation}
the integral being taken over the physical cell of Lebesgue measure $V$, so that
$F(k)$ carries the dimension of $f$ times volume. There is deliberately no
factor $V^{-1}$ in \eqref{eq:fourier-coefficient}.
\end{definition}

\begin{remark}[The pairings carry $V^{-1}$]\label{rem:v-normalisation}
The coefficient \eqref{eq:fourier-coefficient} is the \emph{integral-normalised}
one. Every Fourier pairing written later has the shape
$V^{-1}\sum_k\overline{F(k)}\,G(k)\,\rho(k)$ with a symbol $\rho$: the factor
$V^{-1}$ belongs to the pairing, never to the coefficient. The competing
convention is the orthonormal one attached to the basis of
\cref{ext:fourier-basis}, whose coefficient at $k$ is $V^{-1}F(k)$; the change of
variables between the two is \cref{lem:haar-lebesgue}. Both conventions occur
below, and interchanging them loses exactly one factor of $V=L^2$.
\end{remark}

\subsection{Modes, weights, and which norm enters them}

\begin{definition}[Fourier index set and weights]\label{def:modes}
The \emph{continuum index set} is $\Z^2$; the \emph{lattice index set} is
$I_N:=(\Z/N\Z)^2$, with $N\in\N$ as in \cref{hyp:volume}. For $k\in\Z^2$ put
$\norm{k}_\infty:=\max\{\abs{k_1},\abs{k_2}\}$. The weights of this paper are
functions of this sup norm and of nothing else; the Euclidean norm enters only
the symbols \eqref{eq:lattice-symbol}--\eqref{eq:continuum-symbol}. The weights
are
\begin{equation}\label{eq:weights}
  w_3(k)=
  \begin{cases}
    1, & k=0,\\[2pt]
    \norm{k}_\infty^{-6}, & k\neq0,
  \end{cases}
  \qquad
  w_2(k)=
  \begin{cases}
    0, & k=0,\\[2pt]
    \norm{k}_\infty^{-4}, & k\neq0,
  \end{cases}
\end{equation}
together with the \emph{full} weight $\widetilde w_2$, which agrees with $w_2$
at every $k\neq0$ and takes the value $\widetilde w_2(0)=1$. Thus $w_3$ and
$\widetilde w_2$ are positive at every mode, the origin included, whereas $w_2$
is not.
\end{definition}

The three weights are not interchangeable, and everything below depends on
keeping them apart. The seminorm weight $w_2$ annihilates the zero mode and
enters only where that mode is restored by a separate additive term --- the split
that makes the weighted trace finite
(\cref{thm:uniform-moment,lem:trace-constant}). The full
weight $\widetilde w_2$ is positive everywhere, so the weighted sum it defines
is an honest norm and its sublevel sets are the compact sets of the argument
(\cref{thm:compact-tail}). The weight $w_3$ is positive everywhere for the same
reason as $\widetilde w_2$, and its weighted space carries the measures
(\cref{def:carrier-hs}).

\begin{definition}[The sup norm on the mode lattice]\label{def:weight-supnorm}
For $k\in\Z^2$ the symbol $\norm{k}_\infty$ denotes the sup norm
$\max_i\abs{k_i}$ of \eqref{eq:weights}, and $\abs{k}$ the Euclidean norm
$(\sum_ik_i^2)^{1/2}$. Every weight of this paper is a function of
$\norm{k}_\infty$ alone, every symbol of
\cref{eq:lattice-symbol,eq:continuum-symbol} a function of $\abs{k}$. The two
are equivalent, $\norm{k}_\infty\le\abs{k}\le 2^{1/2}\norm{k}_\infty$ in two
dimensions, and they give the same summability threshold
(\cref{ext:summability}); they are not interchangeable inside an estimate,
because for $k$ with $\abs{k_1}\neq\abs{k_2}$ the two differ and
$\norm{k}_\infty^{-4}\neq\abs{k}^{-4}$.
\end{definition}

\subsection{Negative-Sobolev carriers}

\begin{definition}[Negative-Sobolev sequence spaces]\label{def:carrier-hs}
Let $\mathcal W=(\mathcal W(k))_{k\in\Z^2}$ be a weight with $\mathcal W(k)>0$
for every $k$, and put
\begin{equation}\label{eq:ell2s}
  \ell^2(\Z^2;\mathcal W):=
  \Bigl\{c=(c_k)_{k\in\Z^2}\in\C^{\Z^2}:
  \sum_{k\in\Z^2}\mathcal W(k)\abs{c_k}^2<\infty\Bigr\},
\end{equation}
a Hilbert space for the inner product
$\ip{c}{d}_{\mathcal W}=\sum_k\mathcal W(k)\overline{c_k}d_k$. For $s>0$ let
\begin{equation}\label{eq:full-weight-family}
  \widetilde w_s(k)=
  \begin{cases}
    1, & k=0,\\[2pt]
    \norm{k}_\infty^{-2s}, & k\neq0,
  \end{cases}
  \qquad\text{and set}\qquad
  \HL{s}:=\ell^2(\Z^2;\widetilde w_s).
\end{equation}
The two exponents of \eqref{eq:weights} are the cases $\widetilde w_2$ and
$\widetilde w_3=w_3$ of this family. This is a definition and not a theorem: no
separately defined distribution space is identified here, in keeping with
\cref{sec:scope}. The letter $s$ is the Sobolev order, so the weight of
$\HL{s}$ decays like $\norm{k}_\infty^{-2s}$; where the \emph{decay exponent}
itself is meant below, it is written out rather than denoted by $s$, because the
two differ by a factor of two and the summability threshold is stated in the
decay exponent.
\end{definition}

Two of these spaces occur, with two different jobs, and neither exponent may be
relabelled. The \emph{carrier} is $\HL{3}$: the lattice laws of
\cref{sec:lattice} are Borel probability measures on its completion, and the weak
convergence asserted in \cref{sec:main} is weak convergence there. The
\emph{energy} is $\HL{2}$: the uniform moment of \cref{thm:uniform-moment} is an
$\HL{2}$ quantity. The inclusion $\HL{2}\subseteq\HL{3}$ is
\cref{prop:completion}, and it is the ratio $w_3/\widetilde w_2=\norm{k}_\infty^{-2}$
that tends to zero --- the $\HL{3}$ weight decays \emph{faster}, not the $\HL{2}$ one ---
which is what makes an $\HL{2}$-sublevel set have small $\HL{3}$-tails, hence
relatively compact closure (\cref{thm:compact-tail}).

A point of $\HL{3}$ may be written as its family of Fourier coefficients, with
the weighted norm \eqref{eq:full-weight-family}, or as the family of
\emph{stored coordinates} obtained by applying $\sqrt{\widetilde w_3(k)}$ to
each coefficient; the two differ by the diagonal isometry of
\cref{def:carrier-hminus3}, and \cref{rem:two-families} says which one each
formula below is written in. The deweighting is defined on all of $\HL{3}$
precisely because $\widetilde w_3(k)>0$ at every $k$, the zero mode included
(\cref{eq:weights}) --- which is why the zero mode of the $\HL{3}$ weight is $1$
and not $0$. Computing an $\HL{2}$ norm in coordinates that already carry the
$\HL{3}$ weight multiplies it by a mode-dependent factor, which is the trap
\cref{rem:conventions} warns about.

\subsection{Fixed volume, the two symbols, and a borrowed fact}

\begin{hypothesis}[Fixed physical volume]\label{hyp:volume}
Throughout, $L>0$ and the mass $m>0$ are fixed. The lattice spacing $a>0$ and
the lattice size $N\in\N$ satisfy
\begin{equation}\label{eq:fixed-volume}
  aN=L,
\end{equation}
and the sequence of lattices is indexed by $n\in\N$ through
\begin{equation}\label{eq:index-dictionary}
  N=n+1,\qquad a=a_n=\frac{L}{n+1}\ .
\end{equation}
\end{hypothesis}

\Cref{hyp:volume} also fixes the bookkeeping that the two halves of the paper
index differently. The mathematics below writes $N$ for the lattice size and
passes to $N\to\infty$; the machine development of \cref{sec:lean-overview}
indexes its statements by $n$ and derives $N=n+1$, $a=L/(n+1)$.
\Cref{eq:index-dictionary} is the whole dictionary, with one consequence: the
mesh is never a free parameter, so only $L$, $m$ and $n$ are free below.

Fixed volume buys one thing, and the argument rests on it: a symbol limit.
Because \eqref{eq:fixed-volume} forces $a=L/N\to0$ at fixed $L$, the lattice
symbol below approaches the continuum symbol mode by mode at each fixed $k$, and
it is exactly because the spacing vanishes while $L$ is held fixed that the limit
is not a rescaling of the field. Fixed volume buys no modulus: no rate, no
exponential tail and no quantitative convergence statement is proved here
(\cref{sec:scope}), and the limit is not asserted to be an $L^2$ function.

Two symbols are used from here on, and neither identity is proved here: the
first is the symbol of the operator of
\cref{def:lattice-laplacian,prop:operator}, the second that of the resolvent of
\cref{lem:resolvent}. The lattice configuration space is
$\R^{I_N}$: a configuration assigns a real number to each site $x\in I_N$, and
the site $x$ has the physical representative $ax\in[0,L)^2$. The lattice symbol
and the continuum symbol are
\begin{equation}\label{eq:lattice-symbol}
  \lambda_N(p):=m^2+4a^{-2}\sum_{i=1}^{2}\sin^2\Bigl(\frac{\pi p_i}{N}\Bigr),
  \qquad p\in I_N ,
\end{equation}
\begin{equation}\label{eq:continuum-symbol}
  \lambda(k):=m^2+\Bigl(\frac{2\pi}{L}\Bigr)^{2}\abs{k}^2 ,
  \qquad k\in\Z^2 .
\end{equation}
The expression \eqref{eq:lattice-symbol} is well defined on the residue class
$p_i\in\Z/N\Z$ because $\sin^2(\pi t)$ is $1$-periodic in $t$. Both symbols are
real and at least $m^2>0$, so no denominator formed from them vanishes.

\begin{external}[Torus Fourier basis]\label{ext:fourier-basis}
Let $z=x/L\in[0,1)^2$ and let $\dd z$ be Lebesgue measure on the unit cell, whose
total mass is $1$. The family $z\mapsto e^{2\pi i\,k\cdot z}$, $k\in\Z^2$, is a
complete orthonormal basis of $L^2([0,1)^2,\dd z;\C)$.
\end{external}

\Cref{ext:fourier-basis} is used as a black box, and no other completeness
statement is needed below. It is quoted in this orthonormal form because that is
the form in which tests are paired with the resolvent in \cref{sec:covariance};
the factor that reconciles it with \eqref{eq:fourier-coefficient} is the
$V^{-1}$ of \cref{rem:v-normalisation}, whose proof is \cref{lem:haar-lebesgue}.

\begin{lemma}[Summability of the weights]\label{ext:summability}
In two dimensions both of the following series converge:
\begin{equation}\label{eq:summability}
  \sum_{k\ne0}\norm{k}_\infty^{-4}<\infty,
  \qquad
  \sum_{k\ne0}\norm{k}_\infty^{-6}<\infty .
\end{equation}\end{lemma}

\begin{proof}
The two sums are explicit. The shell $\norm{k}_\infty=q$, $q\ge1$, carries
$(2q+1)^2-(2q-1)^2=8q$ modes, and on that shell $\norm{k}_\infty^{-r}=q^{-r}$, so
each sum is $8\sum_{q\ge1}q^{1-r}$ with $r=4$ respectively $r=6$; such a series
converges exactly when $r>2$. The inequality that matters is therefore $2<4$ and
$2<6$: \emph{the decay exponent of the weight must exceed the dimension}.
\end{proof}

The two values are worth naming, because the constants below carry them:
\[
  \sum_{k\ne0}\norm{k}_\infty^{-4}=8\zeta(3),
  \qquad
  \sum_{k\ne0}\norm{k}_\infty^{-6}=8\zeta(5).
\]

\begin{remark}[The exponents are the two-dimensional ones]\label{rem:dimension}
The threshold just computed is that the decay exponent exceeds $d$, and here
$d=2$. In $d=4$ the decay exponent $4$ of $\widetilde w_2$ would be critical
rather than summable, so the pair $w_2,w_3$ of \eqref{eq:weights} is a
two-dimensional pair; the paper claims nothing about dimension four and offers
no general one. The source of
\cref{sec:covariance} is registered for a general dimension $d$, but this paper
writes it in two dimensions, like everything else here: the test functions are
functions on $\R^2$, and the shell count in the last step of its proof uses
$d=2$. The measure chain therefore uses $d=2$ throughout, with no statement and
no object of this paper claimed in another dimension (\cref{sec:scope}).
\end{remark}

\begin{remark}[Four conventions a reader can get wrong]\label{rem:conventions}
Each of the following is settled above, and each is easy to reverse by accident.
\emph{(i) Integral versus orthonormal normalisation:} \eqref{eq:fourier-coefficient}
carries no $V^{-1}$, the pairing carries it, and the probability-torus form
appears only where a test is paired with the resolvent in \cref{sec:covariance}
(\cref{rem:v-normalisation,lem:haar-lebesgue}). \emph{(ii) Sup norm versus
Euclidean norm:} the weights are stated for $\norm{k}_\infty$ and the symbols
for $\abs{k}$, and the two must not be exchanged inside an estimate
(\cref{def:weight-supnorm}). \emph{(iii) Seminorm versus full weight:} $w_2$
vanishes at the zero mode and is usable only where that mode is supplied
separately, while $\widetilde w_2$ and $w_3$ are positive everywhere and are
what the norm and the compactness argument use
(\cref{def:modes,def:carrier-hs}). \emph{(iv) Carrier versus energy exponent:}
$\HL{3}$ is the carrier of the laws and the topology of weak convergence,
$\HL{2}$ the exponent of the energy and of the uniform moment, and neither may
stand in for the other
(\cref{def:carrier-hs,thm:uniform-moment,sec:tightness}).
\end{remark}
